\documentclass[12pt,a4paper]{article}
\usepackage[T2A]{fontenc}
\usepackage[utf8]{inputenc}
\usepackage[russian]{babel}
\usepackage{geometry}
\usepackage{xcolor}
\usepackage{listings}
\usepackage{needspace}

\geometry{margin=25mm}

\lstdefinestyle{py}{
  language=Python,
  basicstyle=\ttfamily\small,
  keywordstyle=\color{blue!70!black},
  stringstyle=\color{teal!60!black},
  commentstyle=\color{gray!70!black},
  showstringspaces=false,
  frame=single,
  breaklines=true
}
\lstnewenvironment{codeblock}
  {\Needspace{12\baselineskip}\lstset{style=py}}
  {}

\title{Обзор кода приложения Turbine Wear Forecast}
\author{}
\date{}

\begin{document}
\maketitle

\section*{Введение}
Приложение построено как desktop-клиент на PySide6 с хранением данных в SQLite
через SQLAlchemy. Архитектура разделена на доменные сервисы, слой доступа к данным
и интерфейс пользователя.
Такой подход отделяет расчётную часть (прогноз, анализ формул) от UI и упрощает
тестирование отдельных компонентов.

\section*{Общая архитектура и поток данных}
Код организован вокруг трёх основных слоёв: \textbf{UI} управляет сценариями пользователя,
\textbf{доменные сервисы} выполняют вычисления и преобразование данных, а
\textbf{инфраструктура} отвечает за сохранение и извлечение информации из БД.
Переходы между слоями реализованы через сервисы и репозитории, что позволяет
изолировать вычислительную часть от представления и не привязывать её к конкретному виду интерфейса.

Типовой поток данных можно описать так:
\begin{enumerate}
  \item Пользователь создаёт шаблон детали и набор параметров.
  \item Создаётся конкретная деталь на основе шаблона.
  \item Добавляются измерения (значение параметра и часы наработки).
  \item Сервис прогноза строит модель МНК и рассчитывает момент достижения критики.
  \item Виджет графика визуализирует линию прогноза, доверительный интервал и критическую точку.
\end{enumerate}
За счёт такой последовательности легко проследить, как пользовательский ввод
трансформируется в численные расчёты и затем возвращается в UI в виде визуализации и текстовой расшифровки.

\section*{Сценарий пользователя}
Сценарий работы строится вокруг нескольких окон. Сначала создаётся шаблон детали
в разделе ``Шаблоны'', где указывается название и описание. Затем в окне параметров
шаблона вводятся параметры (название, единицы измерения и критическое значение).
После этого на главном экране можно создать деталь, привязав её к выбранному шаблону,
и начать ввод измерений.

При накоплении нескольких измерений пользователь вызывает прогноз, настраивает формулу МНК
и доверительный коэффициент, после чего получает график и текстовое описание.
Отдельно предусмотрен экспорт архива, чтобы сохранить измерения в формате Excel
для внешнего анализа.

\section*{Пример полного цикла работы}
Ниже описан типовой сценарий, показывающий взаимодействие модулей приложения:
\begin{enumerate}
  \item В разделе ``Шаблоны'' создаётся шаблон ``Лопатка ТВД'' с описанием детали.
  \item В окне параметров добавляются параметры ``Толщина кромки'' и ``Высота кромки'',
        задаются единицы измерения и критические значения.
  \item В главном окне создаётся деталь, выбирается шаблон и вводится серийный номер.
  \item Через диалог измерений добавляются точки по каждому параметру с часами наработки.
  \item В окне детали открываются настройки прогноза, вводится формула, например
        линейная модель, и задаются границы коэффициентов.
  \item После нажатия ``Обновить график'' сервис прогноза рассчитывает модель,
        а UI отображает линию МНК, доверительный интервал и критическое время.
  \item Дополнительно пользователь может экспортировать архив и передать данные
        коллегам для внешнего анализа.
\end{enumerate}
Этот сценарий показывает, как данные проходят весь путь от ручного ввода
до аналитического результата и визуализации.

\section*{Структура проекта}
\begin{itemize}
  \item \texttt{src/app/main.py} --- точка входа и запуск UI.
  \item \texttt{src/app/domain/services} --- логика прогнозирования и разбор формул МНК.
  \item \texttt{src/app/infrastructure/db} --- модели БД и репозитории.
  \item \texttt{src/app/ui} --- окна, диалоги и виджеты графика.
  \item \texttt{src/app/ui/windows} --- основные экраны (старт, шаблоны, деталь).
  \item \texttt{src/app/ui/dialogs} --- ввод параметров, измерений, настроек прогноза.
  \item \texttt{src/app/ui/widgets} --- визуализация графика прогноза.
  \item \texttt{src/app/infrastructure/db/session.py} --- подключение к SQLite и сессии.
\end{itemize}
Каждый слой отвечает за свою область: UI только собирает ввод, доменные сервисы
делают вычисления, а репозитории скрывают детали SQL-запросов.

\section*{Библиотеки и фреймворки}
\begin{itemize}
  \item \textbf{PySide6} --- основной UI-фреймворк для окон, диалогов и виджетов.
  \item \textbf{pyqtgraph} --- построение интерактивных графиков и визуализации прогноза.
  \item \textbf{SQLAlchemy} --- ORM и слой доступа к данным SQLite.
  \item \textbf{Alembic} --- управление миграциями схемы БД.
  \item \textbf{Pydantic} --- валидация и типизация структур конфигурации.
  \item \textbf{NumPy} и \textbf{SciPy} --- массивы и численные вычисления, МНК-аппроксимация.
  \item \textbf{scikit-learn} --- вспомогательные метрики качества прогноза.
  \item \textbf{pandas} --- обработка данных при экспорте и подготовке таблиц.
  \item \textbf{SymPy} --- разбор пользовательских формул МНК.
  \item \textbf{openpyxl} --- формирование XLSX при экспорте архива.
\end{itemize}
Связка PySide6 + pyqtgraph даёт богатые возможности визуализации без сложной настройки,
а научный стек (NumPy, SciPy, scikit-learn) обеспечивает корректные и быстрые расчёты.
SQLAlchemy и Pydantic упрощают хранение и валидацию, оставляя код доменной логики
максимально чистым и читаемым.

\section*{Слой данных и инициализация БД}
База данных хранится в файле \texttt{prediction\_wear\_app.db} в корне проекта.
Подключение создаётся один раз и повторно используется через фабрику сессий.
Это позволяет централизованно управлять транзакциями и гарантировать
закрытие соединений даже при исключениях.

\begin{codeblock}
engine = create_engine(DATABASE_URL, echo=False, future=True)
SessionLocal = sessionmaker(bind=engine, autoflush=False, autocommit=False)

@contextmanager
def get_session() -> Session:
    session: Session = SessionLocal()
    try:
        yield session
        session.commit()
    except Exception:
        session.rollback()
        raise
    finally:
        session.close()
\end{codeblock}

Использование контекстного менеджера упрощает код репозиториев:
операции сохранения и чтения автоматически заворачиваются в транзакции,
а откат выполняется при любой ошибке.

\section*{Точка входа}
Приложение инициализирует БД и открывает стартовое окно со списком деталей.
Создание таблиц происходит один раз при запуске, после чего стартует основной
циклический обработчик событий Qt.
\begin{codeblock}
def run() -> int:
    Base.metadata.create_all(bind=engine)

    app = QApplication(sys.argv)
    window = StartWindow()
    window.show()
    return app.exec()
\end{codeblock}

Вызов \texttt{Base.metadata.create\_all} обеспечивает создание таблиц на старте,
что удобно для локального сценария без внешней миграции. Благодаря этому
приложение можно запустить сразу после установки зависимостей, не выполняя
дополнительных команд. В дальнейшем этот подход можно заменить на миграции,
но для учебного проекта он упрощает развёртывание.

\section*{Модель данных}
В БД есть шаблоны деталей, параметры шаблонов, детали, измерения и настройки прогноза.
Ниже показаны ключевые поля параметра и измерения.
Параметр шаблона хранит допустимое критическое значение, а измерение фиксирует
значение на конкретной наработке.
\begin{codeblock}
class TemplateParameter(Base):
    __tablename__ = "template_parameters"
    id: Mapped[int] = mapped_column(Integer, primary_key=True)
    template_id: Mapped[int] = mapped_column(
        ForeignKey("templates.id", ondelete="CASCADE"), nullable=False
    )
    name: Mapped[str] = mapped_column(String(120), nullable=False)
    unit: Mapped[str] = mapped_column(String(32), default="mm", nullable=False)
    critical_value: Mapped[float] = mapped_column(Float, nullable=False)
\end{codeblock}

\begin{codeblock}
class Measurement(Base):
    __tablename__ = "measurements"
    id: Mapped[int] = mapped_column(Integer, primary_key=True)
    part_id: Mapped[int] = mapped_column(ForeignKey("parts.id", ondelete="CASCADE"))
    parameter_id: Mapped[int] = mapped_column(
        ForeignKey("template_parameters.id", ondelete="RESTRICT")
    )
    operating_hours: Mapped[float] = mapped_column(Float, nullable=False)
    value: Mapped[float] = mapped_column(Float, nullable=False)
\end{codeblock}

Модель охватывает следующие сущности:
\begin{itemize}
  \item \textbf{Template} --- шаблон детали, содержит имя и описание.
  \item \textbf{TemplateParameter} --- параметр шаблона с критическим значением.
  \item \textbf{Part} --- конкретная деталь, связанная с шаблоном.
  \item \textbf{Measurement} --- измерение параметра в определённый момент времени.
  \item \textbf{ForecastConfig} --- настройки прогноза для пары ``шаблон + параметр''.
\end{itemize}
Связи настроены через внешние ключи, а уникальные ограничения защищают от
дублирования параметров и измерений.

Конфигурация прогноза хранит формулу и границы параметров, чтобы обеспечить
персональные настройки для каждого параметра шаблона:
\begin{codeblock}
class ForecastConfig(Base):
    __tablename__ = "forecast_configs"
    template_id: Mapped[int] = mapped_column(ForeignKey("templates.id", ondelete="CASCADE"))
    parameter_id: Mapped[int] = mapped_column(
        ForeignKey("template_parameters.id", ondelete="CASCADE")
    )
    lsq_model_formula: Mapped[str] = mapped_column(Text, default="", nullable=False)
    lsq_param_bounds_json: Mapped[str] = mapped_column(Text, default="{}", nullable=False)
    confidence_k: Mapped[float] = mapped_column(Float, default=2.0, nullable=False)
\end{codeblock}

\section*{Слой доступа к данным}
Репозитории инкапсулируют SQL-запросы и обеспечивают CRUD для сущностей. Например,
для измерений задается сортировка по часам наработки и проверка уникальности точки.
Сортировка важна для корректного построения графика и расчёта прогноза во времени.
\begin{codeblock}
class MeasurementRepository:
    def list_by_part_and_parameter(self, part_id: int, parameter_id: int):
        with get_session() as session:
            stmt = (
                select(Measurement)
                .where(
                    Measurement.part_id == part_id,
                    Measurement.parameter_id == parameter_id,
                )
                .order_by(Measurement.operating_hours.asc())
            )
            return session.scalars(stmt).all()
\end{codeblock}

Кроме измерений есть репозитории для шаблонов, параметров, деталей и настроек прогноза.
Каждый репозиторий инкапсулирует операции CRUD, а ошибки ограничений БД
преобразуются в понятные исключения. Например, при создании параметра
с таким же именем для одного шаблона репозиторий возвращает \texttt{ValueError},
который затем показывается пользователю как уведомление.

Отдельно стоит репозиторий \texttt{ForecastConfigRepository}, который поддерживает
операцию ``get\_or\_create\_default''. Она гарантирует наличие дефолтной конфигурации
для выбранного параметра и создаёт её при первом обращении.

\section*{Разбор формулы МНК}
Формула вводится пользователем, разбирается через \texttt{sympy}, из нее
извлекаются параметры и проверяется наличие переменной $x$.
Поддерживаются варианты записи с ``y = ...'' или без левой части, а список
параметров затем используется для подбора коэффициентов.
\begin{codeblock}
def parse_lsq_formula(formula: str) -> ParsedLSQFormula:
    text = formula.strip()
    if "=" in text:
        _, right = text.split("=", 1)
        expr_text = right.strip()
    else:
        expr_text = text
    expr = sp.sympify(expr_text)
    symbol_names = {str(s) for s in expr.free_symbols}
    if "x" not in symbol_names:
        raise ValueError("Формула должна содержать переменную x.")
    parameter_names = sorted(name for name in symbol_names if name != "x")
    return ParsedLSQFormula(expression=expr, parameter_names=parameter_names)
\end{codeblock}

Помимо разбора формулы предусмотрена сериализация границ параметров в JSON.
Это позволяет хранить диапазоны для каждого коэффициента в БД и восстанавливать их
при повторном открытии диалога настроек.
Ниже показано, как границы преобразуются в JSON и обратно:
\begin{codeblock}
def bounds_to_json(bounds: dict[str, tuple[float, float]]) -> str:
    data = {name: [float(low), float(high)] for name, (low, high) in bounds.items()}
    return json.dumps(data, ensure_ascii=False, separators=(",", ":"))

def bounds_from_json(raw: str | None) -> dict[str, tuple[float, float]]:
    if not raw:
        return {}
    data = json.loads(raw)
    # преобразование и валидация словаря
\end{codeblock}
Таким образом, валидация формулы и границ параметров выполняется ещё на этапе ввода,
а некорректные значения не попадают в вычислительную часть.

\section*{Доменные сервисы}
Доменные сервисы сосредоточены в каталоге \texttt{src/app/domain/services} и
представляют ``ядро'' прикладной логики. Они не зависят от UI и не содержат
SQL-запросов, что делает их удобными для тестирования.

Основные сервисы:
\begin{itemize}
  \item \textbf{ForecastService} --- вычисление прогноза, метрик и критического времени.
  \item \textbf{ArchiveService} --- экспорт измерений в ZIP/XLSX.
  \item \textbf{lsq\_formula} --- парсинг формулы и обработка границ параметров.
\end{itemize}
Такое разделение подчёркивает, что бизнес-логика может переиспользоваться
в других интерфейсах (например, в веб-приложении) без переписывания расчётных функций.

\section*{Сервис прогноза}
Сервис прогнозирования выполняет аппроксимацию, вычисляет метрики Bias/RMSE/MAPE
и определяет момент достижения критического значения.
Он получает точки измерений, строит модель МНК и возвращает структуру данных,
которая затем используется и для графика, и для текстового вывода.
Для краткости опущена инициализация вспомогательных переменных, которые
используются в примере ниже: массивы \texttt{x}/\texttt{y}, символы
\texttt{x\_symbol}/\texttt{param\_symbols}, функция \texttt{wrapped\_formula},
начальные параметры \texttt{p0} и границы \texttt{lower}/\texttt{upper}.
\begin{codeblock}
class ForecastService:
    def compute(
        self,
        operating_hours: list[float],
        values: list[float],
        critical_value: float,
        *,
        lsq_model_formula: str = "",
        lsq_param_bounds_json: str | None = None,
        confidence_k: float = 2.0,
    ) -> ForecastResult:
        if len(operating_hours) < 2:
            raise ValueError("Для прогноза нужно минимум 2 измерения.")
        # Разбор формулы и подбор параметров
        parsed = parse_lsq_formula(lsq_model_formula)
        model_fn = sp.lambdify((x_symbol, *param_symbols), parsed.expression, modules="numpy")
        popt, _ = curve_fit(wrapped_formula, x, y, p0=p0, bounds=(lower, upper))
\end{codeblock}

Алгоритм внутри \texttt{compute} можно разбить на несколько шагов:
\begin{enumerate}
  \item Проверка входных данных и сортировка точек по часам наработки.
  \item Разбор формулы и построение функции \texttt{model\_fn} для численного расчёта.
  \item Настройка границ параметров и подбор начального приближения \texttt{p0}.
  \item Аппроксимация \texttt{curve\_fit} и расчёт прогнозных значений.
  \item Вычисление метрик качества и построение доверительного интервала.
  \item Поиск момента достижения критического значения по нескольким кривым.
\end{enumerate}
Возвращаемая структура \texttt{ForecastResult} содержит всё, что нужно UI:
массивы для графика, значения метрик и три оценки критического времени.

Внутри вычислений используются \texttt{np.errstate} и \texttt{nan\_to\_num},
чтобы подавлять переполнения и преобразовывать NaN/Inf в большие конечные числа.
Это защищает оптимизацию от случайных некорректных значений формулы и
позволяет сервису оставаться устойчивым к ошибкам ввода.

Для прогноза формируются две сетки: более плотная \texttt{x\_search} (3000 точек)
используется для поиска критического времени, а итоговая \texttt{x\_grid} (300 точек)
служит для отрисовки графика. Граница прогноза определяется с учётом
параметра \texttt{horizon\_extra}, а при необходимости расширяется до момента
достижения критики. Такая логика гарантирует, что график охватывает интересующий интервал.

\subsection*{Метрики качества}
Внутри сервиса рассчитываются основные метрики прогноза:
\[
\text{Bias} = \frac{1}{n}\sum_{i=1}^{n}(y_i - \hat{y}_i),
\quad
\text{RMSE} = \sqrt{\frac{1}{n}\sum_{i=1}^{n}(y_i - \hat{y}_i)^2}
\]
MAPE вычисляется только при ненулевых фактических значениях:
\[
\text{MAPE} = \frac{100}{n}\sum_{i=1}^{n}\left|\frac{y_i - \hat{y}_i}{y_i}\right|
\]
Эти показатели используются для текста под графиком и для оценки качества модели.

\subsection*{Поиск критического времени}
Метод \texttt{\_critical\_time\_from\_curve} ищет момент пересечения кривой
с критическим значением. Если точного пересечения нет, используется линейная
интерполяция между двумя соседними точками. Это позволяет получить оценку даже
при редкой сетке прогноза и визуально отметить момент критики на графике.

\begin{codeblock}
def _critical_time_from_curve(self, x_grid, y_grid, critical_value):
    diff = y_grid - float(critical_value)
    if np.any(np.isclose(diff, 0.0, atol=1e-9)):
        idx = int(np.where(np.isclose(diff, 0.0, atol=1e-9))[0][0])
        return float(x_grid[idx])
    # иначе ищем пересечение по знаку
    ...
\end{codeblock}
Метод сначала проверяет точное совпадение с критическим значением, затем ищет смену
знака между соседними точками и выполняет интерполяцию. Такое решение достаточно
надёжно для практических сценариев и сохраняет производительность.

\section*{Настройки прогноза}
Параметры формулы и границ настраиваются через отдельный диалог.
Пользователь вводит формулу МНК, после чего список параметров автоматически
подставляется в таблицу границ. Диалог выполняет быструю проверку: если формула
невалидна, таблица очищается, а при подтверждении показывается предупреждение.

\begin{codeblock}
class ForecastConfigDialog(QDialog):
    def _rebuild_bounds_table(self) -> None:
        formula = self.formula_edit.text().strip()
        if not formula:
            self.bounds_table.setRowCount(0)
            return
        parsed = parse_lsq_formula(formula)
        for name in parsed.parameter_names:
            # заполнение таблицы границ
            ...
\end{codeblock}

Хранение границ в JSON даёт возможность сохранять отдельные диапазоны для каждого
параметра и предотвращает некорректный подбор коэффициентов при оптимизации.

\section*{Окна интерфейса}
Стартовое окно управляет деталями и открывает окно конкретной детали.
Через диалог создаётся новая деталь, после чего список сразу обновляется.
\begin{codeblock}
class StartWindow(QMainWindow):
    def on_add_part(self) -> None:
        dialog = PartDialog(self)
        if dialog.exec():
            template_id, name, serial, place, notes = dialog.get_data()
            self.part_repo.create(template_id, name, serial, place, notes)
            self.reload_parts()
\end{codeblock}

Окно ``Шаблоны'' отвечает за справочник типов деталей. Оно позволяет добавлять,
изменять и удалять шаблоны, а также открывать окно параметров, где задаются
единицы измерения и критические значения. Валидация выполняется ещё на уровне UI:
при пустом имени или попытке удалить шаблон с существующими деталями пользователь
получает сообщение об ошибке.

В окне параметров все значения отображаются в таблице, а редактирование проходит
через отдельный диалог:
\begin{codeblock}
class TemplateParametersWindow(QMainWindow):
    def on_add(self) -> None:
        dialog = ParameterDialog(self)
        if dialog.exec():
            name, unit, crit = dialog.get_data()
            self.repo.create(self.template_id, name, unit, crit)
            self.reload()
\end{codeblock}

В окне детали пользователь вводит измерения и строит прогноз с графиком и текстовой расшифровкой.
Метод собирает точки, применяет настройки формулы и передаёт результат в виджет графика.
\begin{codeblock}
def on_refresh_plot(self) -> None:
    param = self._current_parameter()
    x, y = self._collect_points(param.id)
    cfg = self.forecast_config_repo.get_or_create_default(self._part.template_id, param.id)
    result = self.forecast_service.compute(
        operating_hours=x,
        values=y,
        critical_value=param.critical_value,
        lsq_model_formula=cfg.lsq_model_formula,
        lsq_param_bounds_json=cfg.lsq_param_bounds_json,
        confidence_k=cfg.confidence_k,
    )
    self.plot_widget.draw(result, critical_value=param.critical_value)
\end{codeblock}

Помимо построения графика окно детали формирует текстовую расшифровку прогноза.
MAPE переводится в качественную оценку (``отличное'', ``хорошее'' и т.д.), а также
ищется самый ранний критический параметр среди всех параметров шаблона. Это позволяет
быстро увидеть, какая характеристика детали достигнет критики первой.

\begin{codeblock}
def _describe_mape_quality(self, mape: float | None) -> str:
    if mape is None:
        return "Качество модели: MAPE не определён."
    if mape < MAPE_EXCELLENT_THRESHOLD:
        return "Качество модели: отличное."
    ...
\end{codeblock}

\section*{Диалоги ввода данных}
Для большинства операций предусмотрены небольшие диалоги, которые собирают
пользовательский ввод и приводят его к корректному типу данных. Это упрощает
основные окна: вместо проверки каждого поля они получают уже валидированные
значения.

Основные диалоги:
\begin{itemize}
  \item \textbf{PartDialog} --- создание или редактирование детали.
  \item \textbf{MeasurementDialog} --- ввод пары ``часы наработки + значение''.
  \item \textbf{TemplateDialog} --- создание шаблона с названием и описанием.
  \item \textbf{ParameterDialog} --- параметры шаблона и критическое значение.
  \item \textbf{ArchiveDialog} --- выбор папки и имени архива для экспорта.
\end{itemize}

Пример простого диалога измерения:
\begin{codeblock}
class MeasurementDialog(QDialog):
    def get_data(self) -> tuple[float, float]:
        return float(self.hours_spin.value()), float(self.value_spin.value())
\end{codeblock}
Диалоги обеспечивают минимальную валидацию (например, обязательность полей)
и возвращают данные в удобном для репозиториев виде.

\section*{Навигация и взаимодействие окон}
UI построен так, чтобы пользователь всегда оставался в контексте текущей детали.
Главное окно показывает список деталей и набор кнопок справа. Двойной клик по
элементу списка открывает окно детали, а клавиши Enter/Return запускают тот же сценарий.
Окна ``Шаблоны'' и ``Параметры'' открываются как отдельные окна, что позволяет
настраивать справочники, не закрывая основной список.

В окне детали интерфейс разделён на две зоны: слева таблица измерений и график,
справа панель управления и текстовая расшифровка. Такое расположение позволяет
одновременно видеть входные данные и результат прогноза, не переключаясь между вкладками.
Строгие пропорции колонок задаются через \texttt{setColumnStretch}, что обеспечивает
устойчивый макет при изменении размеров окна.

\section*{Табличные представления}
Таблицы используются для отображения измерений и параметров. Они настроены на режим
``только чтение'', чтобы редактирование происходило строго через диалоги и не
создавало неконсистентные данные. Заголовки растягиваются по ширине, а строки
пересчитываются при обновлении списка.

Такой подход упрощает UX: пользователь всегда видит актуальные значения и
не может случайно ввести некорректные данные напрямую в таблицу.

\section*{График прогноза}
Виджет графика визуализирует точки измерений, линию МНК, доверительную зону
и критическое значение.
Используется \texttt{pyqtgraph}: точки измерений рисуются маркерами, прогноз ---
сплошной линией, \texttt{y\_centered} показывает центрированную кривую
(МНК с учётом Bias), а критическая граница --- отдельной горизонтальной линией.
\begin{codeblock}
class ForecastPlotWidget(QWidget):
    def draw(self, result: ForecastResult, critical_value: float) -> None:
        self.plot.plot(result.x_train, result.y_train, pen=None, symbol="o")
        self.plot.plot(result.x_grid, result.y_lsq, pen=pg.mkPen(width=2))
        self.plot.plot(result.x_grid, result.y_centered, pen=pg.mkPen(style=pg.QtCore.Qt.DashLine))
        crit_line = pg.InfiniteLine(pos=critical_value, angle=0, pen=pg.mkPen("r", width=2))
        self.plot.addItem(crit_line)
\end{codeblock}

Для визуализации доверительного интервала используются две кривые
(\texttt{y\_upper} и \texttt{y\_lower}) и полупрозрачная заливка между ними.
Это позволяет пользователю оценить разброс прогнозируемых значений. Кроме того,
на график добавляются вертикальные линии, показывающие интервал критического времени,
если он рассчитан. Такой подход делает прогноз не только числом, но и наглядным
инструментом поддержки решений.

\section*{Экспорт архива}
Сервис архивирования выгружает измерения в ZIP с Excel-файлами по деталям.
Каждая деталь экспортируется в отдельный XLSX, а имена файлов очищаются от
неподдерживаемых символов.
\begin{codeblock}
class ArchiveService:
    def export(self, directory: str, archive_name: str) -> Path:
        dest_dir = Path(directory)
        dest_dir.mkdir(parents=True, exist_ok=True)
        archive_path = dest_dir / f"{archive_name}.zip"
        with zipfile.ZipFile(archive_path, "w", compression=zipfile.ZIP_DEFLATED) as zf:
            for part in self._part_repo.list():
                xlsx_bytes = self._build_part_xlsx(part.id, part.template_id)
                filename = _safe_filename(part.name) + ".xlsx"
                zf.writestr(filename, xlsx_bytes)
        return archive_path
\end{codeblock}

Внутри \texttt{ArchiveService} используется \texttt{openpyxl}, чтобы создавать
книги Excel на лету. Каждый параметр шаблона становится отдельным листом,
а строки содержат часы наработки и измеренное значение. Если у детали нет
параметров, формируется лист ``Нет данных'', чтобы архив был корректным.
Такая схема помогает быстро передать результаты в сторонние аналитические инструменты.

\section*{Обработка ошибок и валидация}
Приложение использует многоуровневую валидацию. На уровне UI проверяются
обязательные поля, диапазоны чисел и наличие выбранных элементов.
Репозитории дополнительно защищают целостность данных через уникальные ограничения,
а возникающие ошибки \texttt{IntegrityError} преобразуются в понятные сообщения.

Сервис прогнозирования выполняет проверку количества точек, совпадения длин списков
и корректности формулы. Ошибки всплывают как исключения, которые перехватываются
в окне детали и отображаются пользователю в виде диалоговых сообщений.

\section*{Сборка, запуск и проверка качества}
Проект рассчитан на Python 3.12+. Для разработки предусмотрены статический анализ,
форматирование и тесты. Типовой набор команд описан в README и включает:
\begin{itemize}
  \item \texttt{ruff check .} --- статический анализ.
  \item \texttt{black --check .} --- проверка форматирования.
  \item \texttt{mypy src} --- строгая проверка типов.
  \item \texttt{pytest} --- запуск тестов.
\end{itemize}
Это обеспечивает единый стиль и раннее обнаружение ошибок.

\section*{Возможности расширения}
Текущая архитектура уже разделяет вычисления, хранение и UI, поэтому
дальнейшие изменения можно внедрять точечно:
\begin{itemize}
  \item добавить новые метрики качества и отображать их в окне детали;
  \item расширить набор экспорта (CSV, JSON) без изменения доменных сервисов;
  \item подключить миграции Alembic для версионирования схемы БД;
  \item интегрировать дополнительные типы графиков (например, трендовые карты).
\end{itemize}
Такой подход делает проект удобной основой для дальнейшего развития.

\section*{Итог}
Код разделён на понятные слои: данные, доменная логика и UI. Основная вычислительная
часть сосредоточена в сервисе прогноза, а визуализация и ввод данных реализованы
в окнах PySide6 и виджете графика.
Такой подход облегчает поддержку: можно менять формулу прогноза или интерфейс
без затрагивания остальных слоёв.

Расширенный обзор показывает, что приложение покрывает полный цикл работы с данными:
от создания справочников до расчёта прогнозов и экспорта результатов. Несмотря на
компактный размер, код содержит все ключевые элементы промышленного проекта:
валидацию, раздельные слои, контролируемый доступ к БД, а также визуализацию.

Благодаря очевидным точкам расширения и аккуратно организованным модулям
приложение может развиваться в сторону более сложной аналитики, дополнительных
форм отчётов или интеграций с внешними системами.

\end{document}
